\documentclass[english,hebrew]{article}
\usepackage[HE8,T1]{fontenc}
\usepackage[utf8]{inputenc}
\setcounter{secnumdepth}{-2}
\setlength{\parindent}{0bp}
\usepackage{calc}
\usepackage{amsmath}
\usepackage{amsthm}

\makeatletter
\@ifundefined{date}{}{\date{}}
\usepackage{graphicx}
\usepackage{geometry}
\newcommand{\limi}{\underset{n\rightarrow\infty}{\lim}}
\newcommand{\limxz}{\underset{x\rightarrow x_{0}}{\lim}}
\newcommand{\limfxz}{\underset{x\rightarrow x_{0}}{\lim}f(x)}
\newcommand{\limfxm}{\underset{x\rightarrow x_{0}^{-}}{\lim}f(x)}
\newcommand{\limfxp}{\underset{x\rightarrow x_{0}^{+}}{\lim}f(x)}
\newcommand{\limxm}{\underset{x\rightarrow x_{0}^{-}}{\lim}}
\newcommand{\limxp}{\underset{x\rightarrow x_{0}^{+}}{\lim}}
\newcommand{\sqrm}{M_{m\times m}(\mathbb{F})}
\newcommand{\seqa}{(a_{n})_{n=1}^{\infty}}
\newcommand{\seqx}{(x_{n})_{n=1}^{\infty}}
\newcommand{\funcdr}{f:D\rightarrow\mathbb{R}}
\newcommand{\der}{\limxz\frac{f(x)-f(x_{0})}{x-x_{0}}}
\usepackage[all]{pdfprivacy}

\makeatother

\usepackage{babel}
\begin{document}
\newgeometry{top=1cm, bottom=2cm, left=1cm, right=1cm}
\title{חשבון אינפיניטסימלי {\beginL 2\endL} - סילבוס הקורס}
\author{נעם לב}

\maketitle
\noindent\fbox{\begin{minipage}[t]{1\columnwidth - 2\fboxsep - 2\fboxrule}%
\begin{center}
האוניברסיטה העברית \textbar{} חשבון אינפיניטסימלי ){\beginL 2\endL}(
- \L{$80132$}
\par\end{center}%
\end{minipage}}
\begin{center}
הסילבוס \textbf{אינו רשמי }ואין להסתמך עליו אלא ככלי עזר.\\
כתבתי אותו לפי סיכומים משנים קודמות\L{$;$} ייתכן שיש בו אי-דיוקים,
וסביר שלא יתאים במלואו לסמסטרים הבאים.\\
למרות זאת, אני מוצא שחלוקת הקורס לנושאים עוזרת לי בהבנתו, ואם גם אחרים
מרגישים כך, אולי הוא ישמש להם נקודת פתיחה נוחה לסידור החומר על פי
פרקים.
\par\end{center}

\medskip{}

\section{פרק 0: פרולוג, הפונקציה המעריכית והאינטגרל הלא מסוים}

\L{$0.1$} פונקציה קדומה והאינטגרל הלא מסוים

\L{$0.2$} אינטגרלים מיידיים ולינאריות

\L{$0.3$} אינטגרציה בחלקים

\L{$0.4$} אינטגרציה בהצבה

\L{$0.5$} הפונקציה המעריכית \L{$\exp$}

\L{$0.6$} הלוגריתם הטבעי \L{$\ln(x)$}

\L{$0.7$} חזקות עם מעריך ממשי ולוגריתם בבסיס כללי

\L{$0.8$} אינטגרלים של חזקות \L{$\sin(x)$} ו-\L{$\cos(x)$}

\L{$0.9$} פונקציות רציונליות ופירוק לשברים חלקיים

\section{פרק 1: פולינום טיילור וכלל לופיטל}

\L{$1.1$} כלל לופיטל למקרה \L{$\frac{0}{0}$}

\L{$1.2$} כלל לופיטל למקרה \L{$\frac{\infty}{\infty}$}

\L{$1.3$} דיפרנציאביליות והקירוב הלינארי

\L{$1.4$} פולינום טיילור \L{$P_{n,f,a}$}

\L{$1.5$} השארית \L{$R_{n,f,a}$} ומשפט היחידות

\L{$1.6$} משפט טיילור: השארית בצורת לגראנז' ובצורת קושי

\L{$1.7$} פולינום טיילור של מכפלה ושל הרכבה

\L{$1.8$} חישוב גבולות, קירובים והערכת שגיאה

\section{פרק 2: האינטגרל המסוים}

\L{$2.1$} חלוקות וסכומי דארבו

\L{$2.2$} אינטגרל עליון ותחתון והגדרת האינטגרביליות

\L{$2.3$} תנאי דארבו לאינטגרביליות ומשפט דארבו

\L{$2.4$} סכומי רימן ושקילות ההגדרות

\L{$2.5$} תכונות האינטגרל המסוים

\L{$2.6$} רציפות במידה שווה, ליפשיציות ומשפט קנטור

\L{$2.7$} אינטגרביליות של פונקציות מונוטוניות ורציפות, שינוי במספר
סופי של נקודות ופונקציות מדרגות

\L{$2.8$} המשפט היסודי, נוסחת ניוטון-לייבניץ הנוסחה היסודית( ונגזרת
של אינטגרל עם גבולות משתנים

\L{$2.9$} משפט הערך הממוצע האינטגרלי, אינטגרציה בחלקים ובהצבה, והשארית
של טיילור בצורה אינטגרלית

\L{$2.10$} שטחים וגבולות של סדרות בעזרת סכומי רימן

\section{פרק 3: האינטגרל הלא אמיתי}

\L{$3.1$} אינטגרל לא אמיתי בתחום לא חסום ושל פונקציה לא חסומה

\L{$3.2$} האינטגרלים של \L{$\frac{1}{x^{p}}$}

\L{$3.3$} מבחני השוואה לפונקציות אי-שליליות

\L{$3.4$} קריטריון קושי

\L{$3.5$} התכנסות בהחלט ובתנאי

\L{$3.6$} אינטגרלים לא אמיתיים מרובים

\L{$3.7$} מבחן דיריכלה

\section{פרק 4: טורי מספרים}

\L{$4.1$} גבול עליון וגבול תחתון של סדרה

\L{$4.2$} טורים, סכומים חלקיים והתכנסות

\L{$4.3$} התנאי ההכרחי להתכנסות וקריטריון קושי

\L{$4.4$} טורים אי-שליליים ומבחני השוואה

\L{$4.5$} מבחן המנה, מבחן השורש ומבחן האינטגרל

\L{$4.6$} משפט לייבניץ, התכנסות בהחלט ובתנאי

\L{$4.7$} הכנסת סוגריים לטור

\L{$4.8$} מבחן דיריכלה ומבחן אבל

\L{$4.9$} שינוי סדר ומשפט רימן

\L{$4.10$} טורים דו-ממדיים ומכפלת טורים

\L{$4.11$} מכפלות אינסופיות

\section{פרק 5: סדרות וטורי פונקציות}

\L{$5.1$} התכנסות נקודתית

\L{$5.2$} התכנסות במידה שווה ותנאי קושי

\L{$5.3$} טורי פונקציות ומבחן \L{$M$} של ויירשטראס

\L{$5.4$} רציפות ואינטגרביליות של פונקציית הגבול

\L{$5.5$} אינטגרציה וגזירה איבר-איבר

\L{$5.6$} טורי חזקות, רדיוס ותחום התכנסות

\L{$5.7$} רציפות, אינטגרציה וגזירה של טורי חזקות

\L{$5.8$} טורי טיילור ופונקציות אנליטיות

\L{$5.9$} משפט הרציפות של אבל

\L{$5.10$} הגדרת \L{$\sin(x)$} ו-\L{$\cos(x)$} בעזרת טורי חזקות
\end{document}
